**Prompt for Claude: Formal Examiner Audit of Chapter 1**

I want you to act as an external MSc examiner in Mathematical Statistics and Actuarial Sciences.

Please perform a structured audit of Chapter 1 of my thesis. Do NOT rewrite it. Instead, evaluate it using the framework below.

---

### 1. Mathematical Rigor (20 marks)

Evaluate:

* Are all definitions precise?
* Is the formulation of martingale measures correct?
* Is the optional decomposition theorem correctly stated?
* Is the distinction between predictable and adapted processes mathematically accurate?
* Are symbols used consistently?
* Are there any subtle stochastic calculus errors?

Provide:

* Mark out of 20.
* Specific technical weaknesses (if any).
* Any notation inconsistencies.

---

### 2. Conceptual Depth (20 marks)

Evaluate:

* Does the chapter demonstrate genuine understanding of incomplete markets?
* Is the pricing interval explained properly?
* Is the minimax structure of American super-hedging clear?
* Is the interpretation of residual cost as the cost of incompleteness convincing?

Provide:

* Mark out of 20.
* Identify any superficial explanations.
* Suggest where intuition could be sharpened.

---

### 3. Theoretical–Computational Bridge (20 marks)

Evaluate:

* Is the neural network clearly interpreted as an estimator of the predictable integrand?
* Is the residual error clearly linked to the adapted cost process?
* Is the difference between Deep BSDE and Deep Hedging technically accurate?
* Is the architecture comparison logically motivated?

Provide:

* Mark out of 20.
* Identify whether any claims are overstated.

---

### 4. Research Gap and Contribution Clarity (20 marks)

Evaluate:

* Is the novelty clearly articulated?
* Is it obvious what has not been done before?
* Is the contribution framed correctly as empirical rather than theoretical?
* Would an examiner understand why this thesis matters?

Provide:

* Mark out of 20.
* Rewrite (in 3–5 sentences) a sharper research gap if needed.

---

### 5. Academic Writing and Structure (20 marks)

Evaluate:

* Is the structure coherent?
* Are transitions smooth?
* Is there unnecessary repetition?
* Is there excessive density?
* Does the tone match MSc-level mathematical writing?

Provide:

* Mark out of 20.
* Identify 3 structural improvements.

---

### 6. Final Verdict

Provide:

* Overall mark out of 100.
* Whether it is:

  * Pass (60–69)
  * Merit (70–79)
  * Distinction (80–89)
  * High Distinction (90+)
* 5 critical improvements that would raise the mark by 5–10%.

Be strict and objective.

---

# 🧠 Why This Works

This forces:

* Structured evaluation.
* Mark allocation.
* Critical weaknesses.
* No fluff rewriting.
* Examiner-style thinking.


